\chapter{如何成为优秀的牌手}

\section{本章导入}

本章将帮助读者先建立对德州扑克的正确认识：它不是依赖短期结果和情绪冲动的游戏，而是一项围绕概率、决策质量、心理控制和风险管理展开的竞技活动。对新手而言，成为优秀牌手的第一步不是记住更多技巧，而是建立长期主义的学习方式。

优秀牌手通常具备三类基础能力：稳定的体能与专注力、成熟的心态管理，以及可复盘的技术框架。本章后续将围绕这三项基础展开，并说明 EV 思维、资金管理和复盘习惯为什么比单手牌的输赢更重要。

\section{核心概念}

本章重点关注长期期望值（EV）、情绪控制、风险管理、资金管理和复盘流程。读者需要理解，短期结果可能受运气影响，但长期收益更接近决策质量的累积结果。

\section{新手常见误区}

新手常把赢钱等同于打得好，把输钱等同于打得差；也容易在连续失利后进入情绪化决策。本节后续将整理这些误区，并引导读者用过程指标替代结果崇拜。

\section{实战思考框架}

面对一手牌，优秀牌手会先问：我的范围是什么，对手的范围是什么，当前行动的目的是什么，风险是否在可承受范围内。这个框架会贯穿全书后续章节。

\section{牌例拆解}

\subsection{牌例一：待补充}

本牌例将用于展示短期结果与长期正确决策之间的区别。

\subsection{牌例二：待补充}

本牌例将用于展示情绪化追损如何破坏原本合理的策略。

\section{本章总结}

本章强调优秀牌手的基础不是神秘直觉，而是长期主义、稳定心态、风险控制和持续复盘。后续章节将在这一基础上进入具体街次决策。

\section{课后训练}

请记录最近三手印象深刻的牌，并尝试分别写下：当时的决策理由、是否受情绪影响、复盘后是否仍然认可该决策。
